\documentclass[12pt]{article}

% ── Packages ──────────────────────────────────────────────────────────────────
\usepackage[margin=1in]{geometry}
\usepackage{hyperref}
\usepackage{booktabs}
\usepackage{array}
\usepackage{parskip}
\usepackage{setspace}
\usepackage{xcolor}
\usepackage{titlesec}
\usepackage{enumitem}
\usepackage{csquotes}
\usepackage[style=numeric,sorting=nyt,backend=biber]{biblatex}

\addbibresource{refs.bib}

% ── Metadata ──────────────────────────────────────────────────────────────────
\hypersetup{
  pdftitle={AI Safety: 3 + 1 --- Governance, Evaluation, Coordination, and the Interpretability Microscope},
  pdfauthor={Mohan Chinnappan},
  pdfkeywords={AI safety, interpretability, mechanistic interpretability,
               AI governance, AI alignment, frontier AI},
  colorlinks=true,
  linkcolor=blue!60!black,
  citecolor=blue!60!black,
  urlcolor=blue!60!black,
}

\titleformat{\section}{\large\bfseries}{}{0em}{}
\titleformat{\subsection}{\normalsize\bfseries}{}{0em}{}

% ──────────────────────────────────────────────────────────────────────────────
\begin{document}

\title{\textbf{AI Safety: 3 + 1}\\
  \large Governance, Evaluation, Coordination,\\
  and the Interpretability Microscope}

\author{Mohan Chinnappan}

\date{September 2026}

\maketitle
\thispagestyle{empty}

% ── Abstract ──────────────────────────────────────────────────────────────────
\begin{abstract}
Dario Amodei has proposed three pillars for safer frontier AI development:
independent evaluation, democratic coordination, and global coordination.
Each addresses a distinct failure mode. Each is necessary. But together they
share a common blind spot: all three operate on observable behavior. None of
them provides a way to look \emph{inside} the system being governed. This
position paper argues that a fourth pillar---\emph{interpretability}---is
required to complete the framework. Using the analogy of a boiler viewport,
we show that governance without internal visibility is governance with a
fundamental instrument gap. We do not argue against the three existing
pillars; we argue that they need a scientific instrument underneath them---one
that transforms behavioral observation into causal understanding. We call this
instrument the microscope, and we argue that building it must be treated as a
first-class safety priority alongside the governance infrastructure it is
meant to inform.
\end{abstract}

\bigskip
\textbf{Keywords:} AI safety, interpretability, mechanistic interpretability,
AI governance, AI alignment, frontier AI, responsible AI.

\bigskip
\noindent\textbf{MSC / ACM codes:} cs.AI, cs.CY, cs.LG

\newpage
\setcounter{page}{1}
\onehalfspacing

% ──────────────────────────────────────────────────────────────────────────────
\section{Introduction}
\label{sec:intro}

The development of frontier artificial intelligence systems has outpaced our
collective ability to understand them. This is not a temporary gap that will
close automatically as capabilities improve. It is a structural problem that
requires deliberate investment to solve.

Anthropic's CEO Dario Amodei articulated one of the clearest public accounts
of what a responsible governance structure for frontier AI development might
look like \cite{amodei2025pace}. His framework identifies three concrete
pillars: independent evaluators with ongoing access to frontier AI companies,
democratic coordination to establish common standards across countries with
shared governance norms, and global coordination to address risks that cross
geopolitical lines.

This paper accepts that framework as necessary. It argues that the framework
is not yet sufficient.

All three pillars operate on the same epistemic layer: observable behavior.
Independent evaluators probe outputs. Democratic coordination sets behavioral
standards. Global coordination manages behavioral risks. None of these
approaches addresses the question that has become increasingly central to the
safety literature: \emph{what is actually happening inside the model?}

We propose a fourth pillar---\emph{interpretability}---as the scientific
instrument that makes the other three enforceable rather than merely
aspirational. The argument proceeds in four steps. Section~\ref{sec:pillars}
describes the three existing pillars. Section~\ref{sec:boiler} introduces the
boiler analogy that motivates the instrument gap. Section~\ref{sec:fourth}
develops the case for interpretability as the fourth pillar.
Section~\ref{sec:together} shows how the four pillars complement rather than
compete with each other. Section~\ref{sec:conclusion} concludes.

% ──────────────────────────────────────────────────────────────────────────────
\section{The Three Pillars}
\label{sec:pillars}

\subsection{Why safety must keep pace with capability}

The starting premise of Amodei's framework---and the one we accept without
reservation---is that capability advancement and safety investment must be
treated as joint obligations. Slowing all AI progress is not the right
response to safety concerns. Ensuring that our ability to understand and
control AI systems advances \emph{alongside} capability is.

The stakes are asymmetric in a way that should shape resource allocation. A
safety failure in a high-stakes deployment---medical, financial, critical
infrastructure---can far exceed the cost of the safety research that would
have prevented it. That asymmetry is not unique to AI; it is the same
asymmetry that motivates safety investment in aviation, nuclear energy, and
pharmaceutical development.

\subsection{Pillar 1: Independent Evaluators}

The first pillar calls for third-party safety evaluators with ongoing access
to frontier AI companies. The key features are:
\begin{itemize}[noitemsep]
  \item Ongoing access, not one-time audits
  \item Independence from the companies being evaluated
  \item Public reporting of significant findings
  \item Authority to flag incidents before they become crises
\end{itemize}

The model here is analogous to financial auditing or pharmaceutical trials:
independent verification of claims that cannot be self-certifying. The
limitation is that such evaluators can only probe external behavior. They
cannot, with current tools, examine the internal mechanisms that produce that
behavior.

\subsection{Pillar 2: Democratic Coordination}

The second pillar calls for common safety standards across countries that
share democratic governance norms. This addresses the collective action
problem in which any single country setting stringent standards faces a
competitive disadvantage unless others follow. Key elements include:
\begin{itemize}[noitemsep]
  \item Shared definitions of high-risk capability
  \item Common deployment thresholds and red lines
  \item Mutual recognition of evaluation findings
  \item Coordinated incident response across borders
\end{itemize}

The limitation is that standards without measurement tools are difficult to
enforce consistently. When the object of measurement is a behavioral surface
rather than a mechanistic ground truth, standards can be gamed or
misapplied.

\subsection{Pillar 3: Global Coordination}

The third pillar extends coordination beyond democratic alliances to address
risks that cross geopolitical lines. This is the most difficult pillar
institutionally, but it is motivated by a straightforward observation: some
risks do not respect borders. An AI-enabled biological threat or a
misaligned system optimizing against human interests would affect populations
regardless of which country's regulatory framework governed its development.
The analogy to nuclear safety communication channels---maintained even between
Cold War adversaries---captures the logic well.

% ──────────────────────────────────────────────────────────────────────────────
\section{The Boiler Analogy: Governing a Flame We Cannot See}
\label{sec:boiler}

Consider the boiler.

A well-designed industrial boiler has a \emph{viewport}: a small window
through which the operator can observe the flame directly. The color of the
flame---blue, yellow, or something anomalous---carries causal information
that no external measurement fully captures. The viewport is not a luxury
feature. It is the difference between understanding what is happening inside
the machine and merely knowing that water is getting hot.

Now imagine operating a powerful boiler with no viewport. The operator can
measure temperature, pressure, and output volume. She can install safety
valves. She can hire inspectors to walk the floor. She can set standards for
acceptable operating ranges. But she cannot see the flame. Her governance of
the system is comprehensive by every external metric---and blind to the
internal process it is meant to govern.

Table~\ref{tab:boiler} contrasts the two situations.

\begin{table}[h]
\centering
\caption{What is observable with and without a viewport.}
\label{tab:boiler}
\begin{tabular}{p{0.46\linewidth}p{0.46\linewidth}}
\toprule
\textbf{Boiler with viewport} & \textbf{Boiler without viewport (today's AI)} \\
\midrule
Measure temperature and pressure & Measure inputs and outputs \\
See the color and behavior of the flame & Evaluate behavioral benchmarks \\
Detect abnormality before failure & Install guardrails and safety valves \\
Understand \emph{why} a reading changed & Hire inspectors to probe externally \\
Intervene precisely, not just reactively & Cannot see what is happening inside \\
\bottomrule
\end{tabular}
\end{table}

Transformer models do not have a conventional program that can be inspected
to explain a particular behavior. Researchers have made significant progress
in mechanistic interpretability and circuit analysis
\cite{elhage2021mathematical,olah2020zoom,conmy2023towards,lieberum2023does}.
But the field is still far from having a complete internal map of a frontier
model. We have increasingly sophisticated ways to measure behavior, evaluate
models, and impose safeguards. None of that is the viewport.

This is what we call the \emph{instrument gap}: the three pillars of
governance are in place---the rules, the inspectors, and the coordination
channels---but the instrument that would give those pillars causal grounding
has not yet been built.

% ──────────────────────────────────────────────────────────────────────────────
\section{The Fourth Pillar: Build the Microscope}
\label{sec:fourth}

Interpretability research aims to build the viewport. If we can understand
which internal circuits produce a particular behavior, we move from observing
\emph{what} an AI system does toward understanding \emph{why} it does it.
That epistemic shift changes the safety equation.

The research program has several components, each addressing a different
dimension of the instrument gap:

\begin{enumerate}[noitemsep]
  \item \textbf{Neural circuits and representations.} Mapping the internal
    structures that encode concepts, relationships, and reasoning paths
    within the model---the basic anatomy before we can understand the
    physiology \cite{elhage2022toy,olah2020zoom}.

  \item \textbf{Mechanistic interpretability.} Understanding how specific
    computations are implemented at the level of weights and activations---not
    just what a model does, but how it does it at the mechanistic level
    \cite{elhage2021mathematical,conmy2023towards}.

  \item \textbf{Causal analysis of model behavior.} Moving from correlation
    (the model outputs $X$ when given input $Y$) to causation (input $Y$
    caused output $X$ because of internal mechanism $Z$)---the difference
    between observation and understanding \cite{geiger2021causal}.

  \item \textbf{Detection of deceptive strategies.} Developing tools that
    can identify when a model is pursuing a strategy that differs from the
    one it appears to be pursuing externally---the capability that would make
    independent evaluation far more powerful \cite{hubinger2019risks}.

  \item \textbf{Tracing representations to actions.} Building methods for
    tracing how internal representations flow through the model and give rise
    to specific outputs---the chain of causation from representation to
    action, made visible.

  \item \textbf{Real-time internal monitoring.} Creating tools that allow
    important internal mechanisms to be observed while the model is running,
    not just after the fact, so that inspectors have something real to inspect.
\end{enumerate}

The stronger and more defensible argument for interpretability is not that it
will predict every rogue action before it happens. It is that greater causal
understanding could substantially improve our ability to detect, explain,
evaluate, and potentially control dangerous behavior---and that this
capability is a prerequisite for the other three pillars to operate at their
full potential.

% ──────────────────────────────────────────────────────────────────────────────
\section{3 + 1: Complementary Layers, Not Competing Approaches}
\label{sec:together}

The four pillars are not in tension. Each addresses a failure mode the others
cannot cover. Table~\ref{tab:fourpillars} summarizes what each pillar
provides, what it cannot do alone, and what it gains from interpretability.

\begin{table}[h]
\centering
\caption{The four-pillar framework: contributions and interdependencies.}
\label{tab:fourpillars}
\renewcommand{\arraystretch}{1.3}
\begin{tabular}{p{0.17\linewidth}p{0.24\linewidth}p{0.24\linewidth}p{0.27\linewidth}}
\toprule
\textbf{Pillar} & \textbf{What it provides} & \textbf{What it cannot do alone} & \textbf{What it gains from interpretability} \\
\midrule
Independent evaluators &
  Third-party verification of safety claims and incident reporting &
  Cannot see inside the model; can only probe external behavior &
  Inspectors gain the ability to examine internal mechanisms directly, not just outputs \\
Democratic coordination &
  Common safety standards across countries with shared governance norms &
  Standards without measurement tools are difficult to enforce consistently &
  Shared interpretability methods give standards a common scientific grounding \\
Global coordination &
  Channels for addressing cross-border risks, including with non-democratic states &
  Difficult to verify compliance when internal model state is opaque &
  Interpretability tools that reveal internal behavior make verification more tractable \\
Interpretability &
  Causal understanding of model internals---the viewport into the machine &
  Technical capability alone does not create accountability structures or global norms &
  The instrument that makes the first three pillars enforceable, not just aspirational \\
\bottomrule
\end{tabular}
\end{table}

The relationship between the pillars can be stated simply:
\begin{itemize}[noitemsep]
  \item \textbf{Governance} tells us what standards we should follow.
  \item \textbf{Independent evaluators} tell us whether those standards are
    actually being followed.
  \item \textbf{International coordination} gives us a way to address risks
    that cross borders.
  \item \textbf{Interpretability} gives us the microscope to understand what
    is happening inside the system---the instrument that makes the first three
    pillars far more powerful, because inspectors with a viewport are
    fundamentally different from inspectors without one.
\end{itemize}

The first three pillars give us the rules, the inspectors, and the
coordination. The fourth pillar gives us the viewport. Without it, even the
best-designed governance system is reading a boiler with no window.

This is not an argument against Amodei's framework. It is an argument that
the framework needs a scientific instrument underneath it---one that turns
governance from behavioral observation into causal understanding.

There is an old engineering principle: if you cannot measure it, you cannot
control it. For increasingly autonomous AI systems, measurement cannot stop
at the input and output surface. We need better instruments for looking
inside. The boiler that gets built right is the one with a viewport.

% ──────────────────────────────────────────────────────────────────────────────
\section{Conclusion}
\label{sec:conclusion}

We have argued that the three-pillar AI safety framework proposed by Dario
Amodei is necessary but incomplete. All three pillars operate on observable
behavior. None provides the causal understanding of model internals that
would allow governance to be grounded in mechanism rather than surface.

We propose interpretability as a fourth pillar: not as an alternative to
governance, but as the scientific instrument that gives governance its teeth.
Independent evaluators with interpretability tools can examine internal
mechanisms, not just outputs. Common standards built on interpretability
methods have a shared scientific grounding. Verification of compliance
becomes more tractable when internal behavior is visible rather than opaque.

The objective is not to slow AI progress. Quite the opposite: we should
accelerate AI research, and accelerate AI safety research even faster. The
instrument has to keep up with what it is being asked to measure.

Build the rules. Hire the inspectors. Coordinate across borders. Build the
microscope.

% ──────────────────────────────────────────────────────────────────────────────
\printbibliography

\end{document}
